\chapter{Azure Synapse Analytics: Dedicated SQL Pools}
\index{Azure Synapse Analytics}\index{dedicated SQL pool}\index{MPP}\index{Data Warehouse Units}\index{dimensional modeling}\index{star schema}%
\begin{objectives}
\begin{itemize}
    \item Explain the massively parallel processing (MPP) architecture of Azure Synapse Analytics dedicated SQL pools, including the control node, compute nodes, and the Data Movement Service.
    \item Relate Data Warehouse Units (DWUs and cDWUs) to performance, concurrency, and hourly cost, and explain pause, resume, and scale operations.
    \item Apply Kimball dimensional modeling: declare the grain, choose among transaction, periodic snapshot, and accumulating snapshot fact tables, and design conformed dimensions.
    \item Implement slowly changing dimensions with SCD Type 1 and SCD Type 2 patterns, including a MERGE-based Type 2 upsert.
    \item Select hash, round-robin, or replicated distribution for a table based on join behavior and data volume.
    \item Bulk-load Parquet data from ADLS Gen2 with the COPY statement and managed identity authentication.
    \item Estimate the monthly cost of a dedicated pool versus a serverless SQL pool for an end-of-month BI workload.
\end{itemize}
\end{objectives}

\begin{crosscloud}
The MPP data warehouse is one of the most portable ideas in analytics engineering. Synapse dedicated SQL pools distribute tables across 60 physical distributions and execute queries with a control node coordinating many compute nodes; the same mental model applies to Amazon Redshift's ra3 clusters and Google BigQuery's slot-based execution.

\vspace{2mm}
{\footnotesize
\begin{tabular}{@{}p{3.0cm}p{2.9cm}p{3.1cm}p{3.2cm}@{}} \\
\toprule
\textbf{Capability} & \textbf{AWS} & \textbf{Azure} & \textbf{GCP} \\
\midrule
MPP warehouse & Amazon Redshift & Synapse dedicated SQL pool & BigQuery (slot-based) \\
Sizing unit & ra3 nodes / RPUs & cDWU (pauseable) & Slots / autoscaling editions \\
Distribution & DISTKEY / DISTSTYLE & Hash / round-robin / replicated & Automatic (no user control) \\
Bulk load from lake & COPY from S3 & COPY INTO from ADLS Gen2 & LOAD DATA from Cloud Storage \\
Serverless SQL over lake & Athena / Redshift Spectrum & Synapse serverless SQL pool & BigQuery external tables \\
Logical warehouse & Redshift views + Spectrum & Serverless views + external tables & BigQuery views \\
\addlinespace
\multicolumn{4}{@{}l@{}}{\textbf{Fabric} warehouse: T-SQL MPP on OneLake; \textbf{Snowflake}: virtual warehouses; \textbf{MongoDB}: Atlas SQL / federated queries (not MPP).} \\
\bottomrule
\end{tabular}}

\vspace{1mm}
The Kimball discipline---grain declarations, conformed dimensions, and slowly changing dimensions---is engine-neutral. A star schema designed well in this chapter would port to Redshift, BigQuery, Snowflake, or a Fabric warehouse with only dialect changes.
\end{crosscloud}

\begin{keyterms}
\textbf{Dedicated SQL pool}: the provisioned MPP data warehouse inside Azure Synapse Analytics, billed per hour while running.\quad
\textbf{DWU / cDWU}: Data Warehouse Unit, the bundled compute-memory-IO sizing metric of a dedicated pool.\quad
\textbf{Distribution}: the placement strategy that decides which of the 60 distributions holds each row.\quad
\textbf{Grain}: the precise business meaning of one row in a fact table.\quad
\textbf{SCD Type 2}: a dimension design that preserves history by closing old rows and inserting new versions.
\end{keyterms}

\section{Week 5 Overview: The Provisioned Warehouse in a Lakehouse World}
Month 2 shifts from storing data to serving it at warehouse scale. Azure Synapse Analytics is Microsoft's umbrella analytics service; inside it, the \textbf{dedicated SQL pool} is the provisioned, MPP data warehouse formerly known as Azure SQL Data Warehouse \cite{microsoft2025synapse}. The pool distributes a table's rows across 60 distributions mapped to compute nodes, and a control node parses queries, builds a distributed execution plan, and coordinates the compute nodes. This architecture is what lets a single query scan hundreds of terabytes---and it is also why placement decisions made at table-creation time dominate query performance later.

Week 5 therefore pairs two disciplines that are often taught separately: the platform mechanics of Synapse dedicated pools (cDWUs, distributions, COPY) and the modeling discipline of Kimball dimensional design (grain, fact types, conformed dimensions, SCD Type 2). Neither is sufficient alone. A perfectly tuned pool running a grain-confused fact table produces fast wrong answers; a beautiful star schema on round-robin distributions produces slow right answers.

\begin{definitionbox}
\textbf{Massively parallel processing (MPP)} is an architecture in which a query is decomposed into parallel steps that execute across many compute nodes, each holding a slice of the data, with a control node assembling partial results into a final answer.
\end{definitionbox}

\begin{notebox}
A dedicated pool bills by the hour whenever it is running, even when idle. Pause or scale down outside working hours, and treat an Azure Monitor cost alert as a mandatory guardrail: an idle DW1000c pool can cost thousands of dollars per month.
\end{notebox}

\section{Monday: MPP Architecture and Dimensional Modeling}
\index{MPP!architecture}\index{control node}\index{compute node}\index{Data Movement Service}\index{Kimball}

\subsection{Control Node, Compute Nodes, and the Data Movement Service}
A dedicated SQL pool has one control node and one or more compute nodes. The control node accepts the T-SQL connection, optimizes the query into distributed operations, and coordinates execution. Compute nodes store distributions of user tables in Azure Storage-backed structures and execute the parallel steps. When a query needs rows that live on different compute nodes---for example, a join between two tables distributed on different keys---the \textbf{Data Movement Service} (DMS) shuffles rows between nodes. Data movement is the single most expensive operation in an MPP system; the craft of warehouse design is mostly the craft of avoiding it.

\begin{definitionbox}
A \textbf{Data Warehouse Unit} bundles CPU, memory, and IO into a single capacity knob. The current generation is expressed as cDWUs (for example DW100c through DW30000c). Scaling up adds compute nodes and raises concurrency limits; pausing stops compute billing while storage continues to bill.
\end{definitionbox}

Sizing decisions should be driven by measured workload behavior, not by the size of the data alone. A 10~TB warehouse serving three analysts differs from a 1~TB warehouse serving three hundred concurrent dashboards. Scale up for load windows, scale down after, and pause when the pool is unused for hours.

\subsection{Declaring the Grain}
Kimball dimensional modeling begins with a contract: \emph{one row in this fact table means exactly one thing} \cite{kimball2013dwtk}. Declaring the grain---``one row per order line,'' ``one row per account per day,'' ``one row per claim status change''---decides what joins are safe, what aggregations are legal, and which dimensions can attach. Skipping the declaration is the most common cause of double-counted revenue in production warehouses, because two fact rows at different grains joined together silently fan out.

\subsection{Fact Table Types}
\begin{itemize}
    \item \textbf{Transaction fact}: one row per business event (a sale, a click, a payment). Highest detail, largest volume.
    \item \textbf{Periodic snapshot fact}: one row per entity per time period (inventory per SKU per day). Answers ``what was the level at point T'' questions.
    \item \textbf{Accumulating snapshot fact}: one row per process instance, updated as milestones complete (order placed, shipped, delivered). Models pipelines with measurable stage durations.
\end{itemize}

\begin{tipbox}
When a business question is ``how long does it take from order to delivery,'' reach for an accumulating snapshot. When it is ``what was on hand on March 31,'' reach for a periodic snapshot. Forcing either question onto a transaction fact produces convoluted SQL and wrong answers.
\end{tipbox}

\subsection{Conformed Dimensions and Slowly Changing Dimensions}
\textbf{Conformed dimensions} are dimensions shared, with identical keys and meanings, across multiple fact tables and data marts. A single \texttt{dim\_date}, \texttt{dim\_customer}, and \texttt{dim\_store} allow a revenue mart and an inventory mart to be combined in one query without reconciliation meetings. Conformance is an organizational achievement as much as a technical one: it requires one team to own the dimension and every other team to consume it.

Slowly changing dimensions handle the fact that descriptive attributes change. \textbf{SCD Type 1} overwrites the old value---correct for typo fixes. \textbf{SCD Type 2} closes the current row (set \texttt{valid\_to}, set \texttt{is\_current = 0}) and inserts a new version, preserving history---required when the business asks ``what did we believe about this customer when the sale happened?''

\begin{lstlisting}[language=SQL,caption={SCD Type 2 upsert on a Synapse dedicated SQL pool: close changed rows, then insert the new versions.},label={lst:ch5-scd2-merge}]
-- SCD Type 2 upsert on a Synapse Dedicated SQL Pool
MERGE INTO dim_customer AS tgt
USING stg_customer AS src
    ON tgt.customer_id = src.customer_id AND tgt.is_current = 1
WHEN MATCHED AND (tgt.email <> src.email OR tgt.city <> src.city) THEN
    UPDATE SET tgt.valid_to = CAST(GETDATE() AS date), tgt.is_current = 0
WHEN NOT MATCHED THEN
    INSERT (customer_id, email, city, valid_from, valid_to, is_current)
    VALUES (src.customer_id, src.email, src.city,
            CAST(GETDATE() AS date), '9999-12-31', 1);

-- insert new versions for the rows just closed by the MERGE
INSERT INTO dim_customer (customer_id, email, city, valid_from, valid_to, is_current)
SELECT src.customer_id, src.email, src.city,
       CAST(GETDATE() AS date), '9999-12-31', 1
FROM stg_customer AS src
JOIN dim_customer AS tgt
    ON tgt.customer_id = src.customer_id AND tgt.is_current = 1
WHERE tgt.email <> src.email OR tgt.city <> src.city;
\end{lstlisting}

\begin{pitfall}
\textbf{Forgetting the second insert.} A MERGE alone closes changed rows but does not insert their new versions. The canonical Synapse SCD2 pattern pairs the MERGE with the follow-up INSERT shown in Listing~\ref{lst:ch5-scd2-merge}. Omitting it silently deletes history.
\end{pitfall}

\section{Tuesday: Distribution Types}
\index{distribution!hash}\index{distribution!round-robin}\index{distribution!replicated}\index{data movement}

\subsection{Hash Distribution}
Hash distribution applies a hash function to one distribution column and assigns each row to one of the 60 distributions. Rows with the same hash value land together, so a join on the distribution key between two hash-distributed tables is \emph{co-located}: each compute node joins its local slices with no data movement. Choose as the distribution key a high-cardinality column that participates in the workload's dominant large joins---\texttt{customer\_id} for a customer-centric warehouse, \texttt{order\_id} for order-line facts. Avoid columns with heavy skew (one value holding 30\% of rows creates one hot distribution that straggles every query) and avoid nullable columns, because all nulls hash to the same distribution.

\subsection{Round-Robin Distribution}
Round-robin assigns rows to distributions in rotation, giving perfectly even storage but no co-location guarantee. Every large join on a round-robin table requires data movement. Use it for staging tables (where the goal is fast load, not fast join) and for small temporary tables where no good distribution key exists.

\subsection{Replicated Distribution}
A replicated table stores a full copy on every compute node. Reads against it never move data, so it is ideal for small dimensions---date, region, status codes---that join to large facts. The cost is paid at write time: every insert, update, and rebuild must write every copy, so replication is wrong for large or frequently updated tables.

\begin{table}[H]
\centering
\small
\begin{tabular}{@{}p{2.6cm}p{4.4cm}p{3.4cm}p{3.2cm}@{}}
\toprule
\textbf{Distribution} & \textbf{Mechanism} & \textbf{Best for} & \textbf{Watch-outs} \\
\midrule
Hash & Hash of one column assigns the row & Large facts and large dimensions with a dominant join key & Skew and nulls concentrate rows \\
Round-robin & Rows rotate across distributions & Staging, loads, no good key & Every large join shuffles data \\
Replicated & Full copy on each compute node & Small conformed dimensions & Writes multiply; rebuilds are heavy \\
\bottomrule
\end{tabular}
\caption{Distribution selection is the highest-leverage physical design decision in a dedicated SQL pool.}
\label{tab:ch5-distributions}
\end{table}

\begin{pitfall}
\textbf{Round-robin production tables.} Leaving a production fact table on the default round-robin distribution makes every large join pay a full shuffle through the Data Movement Service. Benchmark, then hash-distribute on the dominant join key.
\end{pitfall}

\section{Wednesday: Ingestion with COPY and PolyBase}
\index{COPY statement}\index{PolyBase}\index{managed identity!ingestion}

\subsection{COPY: The Modern Bulk Loader}
The \texttt{COPY INTO} statement is the recommended bulk-load path for dedicated SQL pools \cite{microsoft2025synapseCopy}. It reads CSV, Parquet, and ORC directly from ADLS Gen2 or Blob Storage, supports wildcards over folder hierarchies, and loads in parallel across distributions. Crucially, it authenticates with \textbf{managed identity}: grant the pool's system-assigned identity the \texttt{Storage Blob Data Reader} role on the lake container and the load needs no keys, no SAS tokens, and no secrets in code.

\begin{lstlisting}[language=SQL,caption={Bulk-loading Parquet from ADLS Gen2 with COPY and managed identity.},label={lst:ch5-copy}]
COPY INTO stg_sales_raw
FROM 'https://stadlsdemo.dfs.core.windows.net/raw/sales/year=*/month=*/*.parquet'
WITH (
    FILE_TYPE = 'PARQUET',
    CREDENTIAL = (IDENTITY = 'Managed Identity'),
    MAXERRORS = 100,
    ERRORFILE = 'https://stadlsdemo.dfs.core.windows.net/rejects/'
);
\end{lstlisting}

\subsection{PolyBase: The Legacy Path}
PolyBase predates COPY and loads through external tables: an external data source, an external file format, and an external table definition, then \texttt{INSERT INTO ... SELECT} from the external table. PolyBase remains useful when a load needs external-table semantics---for example, repeatedly querying a lake folder as a table from the dedicated pool---but for pure bulk loading, COPY is simpler, faster to authenticate, and the direction Microsoft documentation recommends \cite{microsoft2025synapseDedicatedBest}.

\begin{notebox}
Load into a round-robin staging table, then \texttt{INSERT ... SELECT} or \texttt{CREATE TABLE AS SELECT} (CTAS) into the final hash-distributed, clustered-columnstore table. CTAS is fully parallel and minimally logged; row-by-row inserts into a columnstore produce tiny rowgroups and poor compression.
\end{notebox}

\subsection{The Load-to-Serve Pattern}
A disciplined dedicated-pool pipeline has three stages: land raw files in ADLS Gen2, COPY into round-robin staging tables, then CTAS into hash-distributed, columnstore production tables with the SCD2 MERGE applied for dimensions. This pattern keeps loads fast, keeps joins co-located, and gives a clean point to reject malformed files via \texttt{MAXERRORS} and \texttt{ERRORFILE}.

\section{Thursday Hands-on Project: Ten Million Rows into a Star Schema}
\index{hands-on lab}\index{COPY statement}\index{star schema}
This lab provisions a dedicated SQL pool, loads 10~million synthetic Parquet rows from ADLS Gen2 with COPY, models them as a star schema, and applies the SCD Type 2 MERGE. Work in a sandbox subscription; pause or delete the pool when finished.

\subsection{Provision the Pool}
\begin{lstlisting}[language=bash,caption={Provisioning a Synapse workspace and a small dedicated SQL pool with Azure CLI.},label={lst:ch5-provision}]
az login
az account set --subscription "00000000-0000-0000-0000-000000000000"
$rg = "rg-de-azure-book-week5"
$loc = "eastus"
$ws = "synapse-week5-001"
$pool = "sqlpoolwk5"
$adls = "stadlsweek5001"
$sqladmin = "sqladminuser"
$sqlpwd = Read-Host "SQL admin key" -AsSecureString

az group create --name $rg --location $loc

az storage account create --name $adls --resource-group $rg `
  --location $loc --sku Standard_LRS --kind StorageV2 --hns true

az synapse workspace create --name $ws --resource-group $rg `
  --storage-account $adls --file-system synapsefs `
  --sql-admin-login-user $sqladmin `
  --sql-admin-login-password (ConvertFrom-SecureString $sqlpwd -AsPlainText) `
  --location $loc

az synapse sql pool create --name $pool --workspace-name $ws `
  --resource-group $rg --performance-level DW100c

az synapse workspace firewall-rule create --name allowClient `
  --workspace-name $ws --resource-group $rg `
  --start-ip-address "<your-ip>" --end-ip-address "<your-ip>"
\end{lstlisting}

Grant the workspace's system-assigned managed identity \texttt{Storage Blob Data Contributor} on the storage account so COPY can read the lake without secrets.

\subsection{Generate and Stage Synthetic Data}
Use the Python \texttt{faker} and \texttt{pyarrow} packages to write 10~million order rows as partitioned Parquet files (\texttt{year=/month=/}), then upload with \texttt{az storage fs file upload} or AzCopy. Ten million rows is large enough that distribution and columnstore choices are observable but small enough to load on DW100c in minutes.

\subsection{Model the Star and Apply SCD2}
Declare the grain first: \emph{one row per order line}. Create \texttt{fact\_sales} hash-distributed on \texttt{customer\_id} with a clustered columnstore index; create \texttt{dim\_date} as replicated; create \texttt{dim\_customer} hash-distributed on \texttt{customer\_id}. Load staging with COPY (Listing~\ref{lst:ch5-copy}), CTAS into production tables, then run the SCD2 MERGE from Listing~\ref{lst:ch5-scd2-merge} against a changed-customer staging batch and verify history with \texttt{SELECT ... WHERE is\_current = 1}.

\subsection{Verify and Pause}
Confirm row counts match between lake files and staging, confirm the columnstore rowgroups are compressed (\texttt{sys.pdw\_nodes\_column\_store\_row\_groups}), and confirm a sample customer shows two Type 2 versions. Then pause the pool:
\begin{lstlisting}[language=bash,caption={Pausing the pool to stop compute billing.},label={lst:ch5-pause}]
az synapse sql pool pause --name $pool --workspace-name $ws --resource-group $rg
\end{lstlisting}

\section{Friday Use Case: End-of-Month BI on a Budget}
\index{use case}\index{cost optimization!Synapse}
A financial services firm runs heavy BI only during the last five business days of each month: roughly 120 concurrency-hours of report load per month on 8~TB of curated data. The architecture question: dedicated pool, serverless pool, or a hybrid?

\subsection{Cost Model for Dedicated}
A dedicated pool bills per hour while running. A DW500c pool running 12 hours per day for five days is 60 hours---cheap. The same pool left running all month is roughly 720 hours---an order of magnitude more expensive with no additional value. The dedicated answer is viable only with disciplined automation: scale up at month-end, pause immediately after. Missed pauses, not the rate card, are the risk.

\subsection{Cost Model for Serverless}
The serverless SQL pool bills per TB scanned. If the month-end dashboards are well-partitioned Parquet with aggressive pruning and scan 2~TB per report cycle, the serverless bill can be a small fraction of a dedicated pool's idle spend---and there is nothing to pause. But cost scales with bytes scanned: a dashboard that scans CSV, or that re-scans a full 8~TB folder because queries lack partition filters, inverts the comparison immediately \cite{microsoft2025synapseServerless}.

\subsection{Recommendation}
For this workload---predictable calendar, high concurrency in bursts, curated Parquet---recommend a \textbf{hybrid}: serverless SQL views over the gold lake zone for exploration and lightweight reporting, with a DW500c dedicated pool scaled in for the month-end crunch when concurrency and latency SLAs are tight. The decision driver is not the rate card; it is the ratio of busy hours to idle hours and the discipline of the pause automation.

\begin{table}[H]
\centering
\small
\begin{tabular}{@{}p{3.2cm}p{3.1cm}p{3.5cm}p{3.8cm}@{}}
\toprule
\textbf{Option} & \textbf{Billing model} & \textbf{Best for} & \textbf{Watch-outs} \\
\midrule
Synapse dedicated SQL pool & cDWU per hour; pauseable & Predictable, high-concurrency warehouse workloads & Idle pools still bill; pause or resize off-hours \\
Synapse serverless SQL & Per TB of data scanned & Ad-hoc exploration, logical warehouse over the lake & Cost scales with bytes scanned; prune partitions, use Parquet \\
Fabric capacity (F SKU) & Reserved capacity units, per-second billing, pauseable & Unified lakehouse plus Power BI on one capacity & Burst-then-throttle when oversubscribed; size the SKU \\
Azure SQL Database & vCore/hour or DTU; serverless option & OLTP systems and small departmental marts & Not MPP; scale ceiling below Synapse for large analytics \\
\bottomrule
\end{tabular}
\caption{Azure compute cost comparison for warehousing workloads.}
\label{tab:ch5-cost-comparison}
\end{table}

\section{Architectural Decision Record Template}
\index{architectural decision record}
Write an ADR for every dedicated-pool deployment: the grain declaration for each fact table; the distribution key per table with measured skew evidence; the cDWU sizing and the scale/pause schedule; the COPY authentication model; and the cost ceiling with the alert threshold that enforces it. An ADR that names the rejected alternatives---round-robin for the fact table, SAS-token loads, 24/7 running---is more useful in an incident than one that only describes the chosen design.

\section{Operational Deep Dive: Concurrency, Skew, and Rowgroup Health}
\index{concurrency}\index{data skew}\index{rowgroup}
Three metrics predict dedicated-pool health before users complain. First, \textbf{concurrency slot pressure}: each cDWU level grants a fixed number of concurrency slots; when queued queries grow, either scale up or reduce concurrent load rather than letting queues lengthen silently. Second, \textbf{distribution skew}: query \texttt{sys.dm\_pdw\_request\_steps} and the distribution-level row counts; if one distribution holds multiples of the median, the hash key is wrong. Third, \textbf{rowgroup quality}: columnstore rowgroups should sit near one million rows; frequent small loads or updates fragment them, and \texttt{ALTER INDEX ... REORGANIZE} or a CTAS rebuild restores compression.

\begin{tipbox}
Schedule a weekly report of the ten largest tables with their distribution type, row count per distribution, and average rowgroup size. Skew and fragmentation are slow leaks; a dashboard makes them visible before the quarter-end query that exposes them.
\end{tipbox}

\section{Troubleshooting Patterns}
\index{troubleshooting!Synapse dedicated pool}
Slow joins usually mean data movement: check the query plan for \texttt{ShuffleMoveOperation} and re-examine distribution keys. COPY failures with authorization errors mean the managed identity lacks a data-plane RBAC role---Owner on the resource is not enough. \texttt{MAXERRORS} rejections that cluster on one folder usually mean a schema drift in one source file. Sudden cost spikes almost always mean a pool that was never paused; check the pause automation before checking the rate card.

\section{Week 5 Lab Rubric}
\index{rubric}
\begin{table}[H]
\centering
\small
\begin{tabular}{@{}p{3.2cm}p{5.0cm}p{5.0cm}@{}}
\toprule
\textbf{Criterion} & \textbf{Meets expectation} & \textbf{Needs improvement} \\
\midrule
Modeling & Grain declared; fact type chosen deliberately; SCD2 verified with a two-version customer & Grain implied; SCD2 MERGE run but history never queried \\
Distribution & Hash keys justified with join evidence; replicated small dims & Everything left round-robin without measurement \\
Ingestion & COPY with managed identity; rejects captured via \texttt{ERRORFILE} & Keys or SAS tokens in code; silent load errors \\
Cost & Pool paused or deleted; cost estimate documented & Pool left running; no cost evidence \\
\bottomrule
\end{tabular}
\caption{Week 5 rubric for the dedicated-pool deliverables.}
\label{tab:ch5-rubric}
\end{table}

\section{Interview Q\&A}
\subsection{Concepts and Trade-offs}
\begin{enumerate}
    \item \textbf{Q: Why must you declare the grain before designing a fact table?}\\
    A: The grain defines what one row means; without it, joins fan out and aggregates double-count silently.
    \item \textbf{Q: Which fact type models a process with milestones such as order, ship, deliver?}\\
    A: An accumulating snapshot---one row per process instance, updated as each milestone completes.
    \item \textbf{Q: When is SCD Type 2 required instead of Type 1?}\\
    A: When the business must reconstruct what was known at a past point in time, such as the customer's segment when an order was placed.
    \item \textbf{Q: When does a serverless SQL pool beat a dedicated pool on cost?}\\
    A: When queries are infrequent or bursty, data is columnar and partition-pruned, and the bytes scanned per month cost less than keeping a provisioned pool running.
\end{enumerate}

\section{Further Reading}
Start with the Azure Synapse Analytics overview \cite{microsoft2025synapse} and the dedicated SQL pool best practices \cite{microsoft2025synapseDedicatedBest}. Study the COPY statement reference \cite{microsoft2025synapseCopy} before any bulk-load design. For dimensional modeling, Kimball and Ross remain the canonical text \cite{kimball2013dwtk}. Compare with the serverless SQL pool documentation \cite{microsoft2025synapseServerless} ahead of Chapter~7, and connect sizing and pause discipline to the Well-Architected cost pillar \cite{microsoft2025wellarchitected}.

\section*{Key Takeaways}
\begin{itemize}
    \item A dedicated SQL pool is an MPP warehouse: one control node coordinates compute nodes, and the Data Movement Service shuffles rows when distributions do not align.
    \item cDWUs bundle compute, memory, and IO; billing is hourly while running, so pause and scale discipline dominates the cost story.
    \item Declare the grain before building any fact table; choose transaction, periodic snapshot, or accumulating snapshot deliberately.
    \item SCD Type 2 in Synapse is a MERGE to close old rows \emph{plus} a second INSERT for the new versions.
    \item Hash-distribute large tables on the dominant join key, replicate small dimensions, and round-robin only staging.
    \item COPY with managed identity is the modern bulk loader; PolyBase external tables remain for external-table semantics.
\end{itemize}

\section{Exercises}
\begin{enumerate}
    \item \textbf{(Conceptual)} Explain the roles of the control node, compute nodes, and the Data Movement Service during a hash-mismatched join.
    \item \textbf{(Conceptual)} A fact table stores one row per account balance per night. Which fact type is this, and what questions does it answer well?
    \item \textbf{(Conceptual)} Why is \texttt{customer\_segment} a poor hash distribution key for \texttt{dim\_customer}?
    \item \textbf{(Hands-on)} Run the Thursday lab, then compare load time of COPY into a round-robin staging table versus directly into a hash-distributed table. Explain the difference.
    \item \textbf{(Hands-on)} Modify the SCD2 MERGE to also track \texttt{loyalty\_tier}. Demonstrate a three-version customer.
    \item \textbf{(Hands-on)} Write a query against \texttt{sys.pdw\_nodes\_column\_store\_row\_groups} that reports compressed versus open rowgroups per table.
    \item \textbf{(Design)} A warehouse has 400 concurrent dashboard users from 08:00--18:00 weekdays only. Propose a scale and pause schedule and estimate the monthly cDWU-hours.
    \item \textbf{(Design)} Given \texttt{fact\_orders} (hash on \texttt{order\_id}) and \texttt{fact\_shipments} (hash on \texttt{shipment\_id}), explain the data movement cost of joining them and propose a redesign.
    \item \textbf{(Design)} Write an ADR choosing DW1000c over DW500c for a quarter-end close. Include the evidence you would collect to defend the choice.
    \item \textbf{(Interview)} Prepare a two-minute explanation of why ``the pool is big enough'' is not an answer to a slow-query complaint.
\end{enumerate}

\section{Capstone Project: Star Schema Warehouse with Managed SCD2 History}\label{sec:ch5-capstone}
\index{capstone project}\index{star schema}\index{SCD Type 2}

\begin{capstonebox}
\textbf{Mission:} Build a small but production-shaped dedicated SQL pool warehouse: load 10~million rows from ADLS Gen2 with COPY and managed identity, model a declared-grain star schema with hash-distributed facts, replicated dimensions, and SCD Type 2 customer history, and demonstrate a measured scale/pause cost plan.

\textbf{Estimated time:} 4--6 hours.

\textbf{What you will practice:}
\begin{itemize}
    \item Provisioning a Synapse workspace and dedicated pool with Azure CLI and firewall-scoped access.
    \item Bulk-loading partitioned Parquet with COPY using the workspace managed identity.
    \item Declaring grain, choosing fact types, and implementing conformed dimensions.
    \item Running the two-step SCD2 MERGE and verifying historical versions.
    \item Hashing on the dominant join key and replicating small dimensions, with skew evidence.
\end{itemize}
\end{capstonebox}

\subsection*{Scenario}
A retail analytics team needs a warehouse over five years of order-line history. Analysts join sales to customers, dates, and stores thousands of times per day; the customer dimension changes weekly as segments and addresses are corrected. Finance requires that any historical report can be re-run as of a past date, so customer history must be preserved, not overwritten.

\subsection*{Requirements}
\begin{itemize}
    \item \textbf{Grain contract:} \texttt{fact\_sales} is one row per order line, stated in the README and enforced by a uniqueness check.
    \item \textbf{Distribution:} \texttt{fact\_sales} and \texttt{dim\_customer} hash on \texttt{customer\_id}; \texttt{dim\_date} and \texttt{dim\_store} replicated; staging round-robin.
    \item \textbf{Ingestion:} COPY with \texttt{CREDENTIAL = (IDENTITY = 'Managed Identity')}; malformed files rejected to an error folder, never silently skipped.
    \item \textbf{History:} SCD Type 2 on \texttt{dim\_customer} with \texttt{valid\_from}, \texttt{valid\_to}, and \texttt{is\_current}; an as-of query demonstrates reconstruction of a past state.
    \item \textbf{Cost:} documented scale/pause schedule; pool paused at submission time; a cost estimate with the DWU-hours assumption.
\end{itemize}

\subsection*{Reference Architecture}
\begin{figure}[H]
    \centering
    \resizebox{\textwidth}{!}{%
    \begin{tikzpicture}[node distance=1.4cm]
        \node[stage] (gen) {Synthetic\\generator};
        \node[store, right=1.6cm of gen] (lake) {ADLS Gen2\\raw Parquet};
        \node[stage, right=1.6cm of lake] (copy) {COPY INTO\\(managed identity)};
        \node[store, right=1.6cm of copy] (stg) {Round-robin\\staging};
        \node[store, right=1.6cm of stg] (fact) {fact\_sales\\hash + CCI};
        \node[store, above=1.1cm of fact] (dims) {Replicated\\dims};
        \node[store, below=1.1cm of fact] (dimc) {dim\_customer\\SCD2};
        \node[stage, right=1.5cm of fact] (bi) {Analyst\\queries};
        \node[doc, below=1.1cm of stg] (cost) {Pause/scale\\schedule};
        \draw[arrow] (gen) -- (lake);
        \draw[arrow] (lake) -- (copy);
        \draw[arrow] (copy) -- (stg);
        \draw[arrow] (stg) -- (fact);
        \draw[arrow] (dims) -- (fact);
        \draw[arrow] (dimc) -- (fact);
        \draw[arrow] (fact) -- (bi);
        \draw[arrow] (cost) -- (stg);
    \end{tikzpicture}%
    }
    \caption{Capstone architecture: lake-staged Parquet, COPY with managed identity, round-robin staging, hash-distributed columnstore facts, and replicated dimensions.}
    \label{fig:ch5-capstone-arch}
\end{figure}

\subsection*{Implementation Steps}
\begin{enumerate}
    \item Provision the workspace and pool from Listing~\ref{lst:ch5-provision}; assign the managed identity \texttt{Storage Blob Data Contributor}.
    \item Generate 10~M order-line Parquet files partitioned by \texttt{year/month}; upload to \texttt{raw/sales/}.
    \item COPY into \texttt{stg\_sales\_raw}; confirm \texttt{MAXERRORS} rejects land in \texttt{rejects/}.
    \item Create the star schema; CTAS staging into \texttt{fact\_sales}; load dimensions.
    \item Run the SCD2 MERGE (Listing~\ref{lst:ch5-scd2-merge}) against a batch with changed emails; verify two versions.
    \item Benchmark a star join before and after fixing distributions; record the timings.
    \item Pause the pool; record the monthly cost estimate.
\end{enumerate}

\subsection*{Deliverables}
\begin{itemize}
    \item SQL scripts for staging, star schema, COPY, CTAS, and the SCD2 MERGE.
    \item A one-page grain declaration and distribution rationale per table.
    \item Query output proving two SCD2 versions and an as-of reconstruction.
    \item Before/after join timings and a monthly cost estimate with pause evidence.
\end{itemize}

\subsection*{Acceptance Criteria}
\begin{itemize}
    \item Row counts reconcile between lake files, staging, and the fact table.
    \item No credentials appear in any script; all loads use managed identity.
    \item \texttt{is\_current = 1} returns exactly one row per customer.
    \item The final pool state is paused or deleted.
\end{itemize}

\subsection*{Stretch Goals}
\begin{itemize}
    \item Add a periodic snapshot fact for daily inventory and justify the grain in the README.
    \item Measure distribution skew per table and document the fix if any distribution exceeds twice the median.
    \item Re-run the COPY benchmark at DW200c and DW400c and plot load time versus cDWU.
\end{itemize}
